\begin{definition}\label{def:norma-general}
   Una \emph{norma} en $\mathbb{R}^n$ es una función escalar, la cual denotaremos por $$\|\cdot\|: \mathbb{R}^n \to \mathbb{R},$$ que satisface las siguientes condiciones, para todo $\mathbf{x}, \mathbf{y} \in \mathbb{R}^n$ y $\lambda \in \mathbb{R}$:
    \begin{enumerate}
        \item Positividad: $\|\mathbf{x}\| \geq 0\;\land\;\|\mathbf{x}\| = 0 \iff \mathbf{x} = \boldsymbol{0}$.
        \item Homogeneidad: $\|\lambda \mathbf{x}\| = |\lambda|\|\mathbf{x}\|$.
        \item Desigualdad triangular: $\|\mathbf{x} + \mathbf{y}\| \leq \|\mathbf{x}\| + \|\mathbf{y}\|$.
    \end{enumerate}
\end{definition}

\begin{obs}
  Existen infinitas funciones escalares que cumplen con las condiciones reci\'en mencionadas, mostraremos algunas de ellas a continuaci\'on.
\end{obs}

\begin{definition}\label{def:norm1}   Sean $\mathbf{x}=(x_1,\ldots,x_n) \in\Rn{n}$. Se define la \emph{norma 1},  la cual denotaremos por $\| \; \|_1,$  a 
    \[
        \|\mathbf{x}\|_1 = \sum_{i=1}^n |x_i|.
    \] 
 \end{definition}

\begin{definition}\label{def:norm2}
  Sean $\mathbf{x}=(x_1,\ldots,x_n) \in\Rn{n}$.  Se define la \emph{norma 2} o \emph{norma euclidiana}, la cual notaremos por  $\|\;\|_2$, a
    \[
        \|\mathbf{x}\|_2 = \sqrt{ \sum_{i=1}^n x_i^2 }.
    \]
  \end{definition}

\begin{tarea} Probar que  $\| \; \|_1$ y  $\| \; \|_2$ son efectivamente normas, es decir, que cumplen con  lo pedido en (\ref{def:norma-general}). 
\end{tarea}

\begin{definition}
    Sean $\mathbf{x},\;\mathbf{y}\in\Rn{n}$. Se define la  \emph{ función distancia} entre $\mathbf{x}\text{ e }\mathbf{y}$ como la funci\'on $d:\Rn{n}\times\Rn{n}\to\R$ tal que
    \[
        d(\mathbf{x},\mathbf{y})=\|\mathbf{x}-\mathbf{y}\|,
    \]
    donde $\|\cdot\|$ es  cualquier norma de $\Rn{n}$.
\end{definition}
